\subsection{nml\_testcase\_channel}
\label{SEC:nml_testcase_channel}
\inp*[nmlTestcaseChannel]{}[nml\_testcase\_channel][]

This \hyperref[SEC:User_Input]{namelist block} is used to specify inputs for the
rectangular channel flow test case
\\% Prevent next word from entering right margin
($\inp[nmlInitialization::Initialization]{Initialization}=2$) or the
annular-pipe channel flow test case ($\inptext{Initialization}=25$).
%
The inputs that start with \quoterm{Grid\_} are specific to running one of these
cases using the internal grid generator.

\begin{nmlblocklst}[
,label={LST:nml_testcase_channel}
,caption={Namelist nml\_testcase\_channel.}
]
&nml_testcase_channel
  Grid_Annulus_Angle  =   5.0    ! [deg]
  Grid_Ave_Xplus      = -30.0
  Grid_Ave_Yplus      = -15.0
  Grid_Ave_Zplus      = -15.0
  Grid_Yplus_Wall     =   2.0
  Grid_Yplus_Of_One   = 8.42e-6  ! [m]
  Body_Force          =   0.0    ! [N]
  Nondim_Output_Coord = F
/
\end{nmlblocklst}

%==========================================================================

%##%\FINALFIX{\newpage}
\begin{nmldescitem}
Grid_Annulus_Angle = 5.0  ! [deg]
\end{nmldescitem}
\inp*[nmlTestcaseChannel::GridAnnulusAngle]{}[Grid\_Annulus\_Angle][]

\begin{nmldescbody}
%
This input is used when running the annular-pipe channel flow case and
specifies the angle (in degrees) of the circumferential extent of the grid
domain.
%
The input value will be limited to an angle between 0.1 and 90 degrees.
%
The radial distance between the inner and outer walls of the annular pipe
will be set to the channel height used for the rectangular channel flow case.
%
The arc length of the annular pipe \emph{outer wall} will be set to the span used for
the regular channel flow case.
%
From the arc length and circumferential angle, the outer radius of the annulus
will be computed.
%
This means the amount of curvature for the annular pipe will vary depending on
the input \inptext{Grid\_Annulus\_Angle}, with a maximum value of 90 resulting in a
highly curved annular pipe and a minimum value of 0.1 have little-to-no
curvature.
%
\end{nmldescbody}

%==========================================================================

\begin{nmldescitem}
Grid_Ave_Xplus = -30.0
\end{nmldescitem}
\inp*[nmlTestcaseChannel::GridAveXplus]{}[Grid\_Ave\_Xplus][]

\begin{nmldescbody}
When running either of the channel flow problems using the internal grid generator,
this input is used to determine the average grid spacing in the
axial ($x$) direction as follows:
\\
\begin{tabular}{ @{\qquad } r @{ \quad } p{0.75\linewidth} @{} }
  $<0$ & The abs(input value) specifies the average \yp[x] degree of freedom
         spacing. This means the average \yp[x] cell spacing is given by
         $\left|\inptext{Grid\_Ave\_Xplus}\right|\times(\p[+1]{S})$.\\
  $>0$ & The input value specifies the average \yp[x] cell spacing.
\end{tabular}
%
\end{nmldescbody}

%==========================================================================

\begin{nmldescitem}
Grid_Ave_Yplus = -15.0
\end{nmldescitem}
\inp*[nmlTestcaseChannel::GridAveYplus]{}[Grid\_Ave\_Yplus][]

\begin{nmldescbody}
When running either of the channel flow problems using the internal grid generator,
this input is used to determine the average grid spacing in the vertical
($y$) direction specifically at the center-height of the channel.
%
Note, for the annular-pipe channel flow case, this is applied to the radial
spacing.
\\
\begin{tabular}{ @{\qquad } r @{ \quad } p{0.75\linewidth} @{} }
  $<0$ & The abs(input value) specifies the average \yp{} degree of freedom
         spacing at the center of the channel. This means the average \yp{}
         cell spacing at the center of the channel is given by
         $\left|\inptext{Grid\_Ave\_Yplus}\right|\times(\p[+1]{S})$.\\
  $>0$ & The input value specifies the average \yp{} cell spacing at the
         center of the channel.
\end{tabular}

\noindent\NOTE
This input will be used to determine the \yp{} spacing for the cells
at the middle of the channel height, and the input
\inp[nmlTestcaseChannel::GridYplusWall]{Grid\_Yplus\_Wall}
will be used to determine the \yp{} spacing for the cells on channel walls.
%
A tanh spacing will then be used to get the \yp{} cell spacings for
all the cells in between.
\end{nmldescbody}

%==========================================================================

\begin{nmldescitem}
Grid_Ave_Zplus = -15.0
\end{nmldescitem}
\inp*[nmlTestcaseChannel::GridAveZplus]{}[Grid\_Ave\_Zplus][]

\begin{nmldescbody}
When running either of the channel flow problems using the internal grid generator,
this input is used to determine the average grid spacing in the
spanwise ($z$) direction.
%
Note, for the annular-pipe channel flow case, this is applied to the
circumferential spacing.
\\
\begin{tabular}{ @{\qquad } r @{ \quad } p{0.80\linewidth} @{} }
  $<0$ & The abs(input value) specifies the average \yp[z] degree of freedom
         spacing. This means the average \yp[z] cell spacing is given by
         $\left|\inptext{Grid\_Ave\_Zplus}\right|\times(\p[+1]{S})$.\\
  $>0$ & The input value specifies the average \yp[z] cell spacing.
\end{tabular}
\end{nmldescbody}

%==========================================================================

\begin{nmldescitem}
Grid_Yplus_Wall = 2.0
\end{nmldescitem}
\inp*[nmlTestcaseChannel::GridYplusWall]{}[Grid\_Yplus\_Wall][]

\begin{nmldescbody}
When running either of the channel flow problems using the internal grid generator,
this input is used to determine the height of the first cell off the wall.
%
Note, for the annular-pipe channel flow case, this is applied to the radial wall
spacing.
\\
\begin{tabular}{ @{\qquad } r @{ \quad } p{0.75\linewidth} @{} }
  $<0$ & The location of \yp{}=1 is located at the
         $\mathtt{nint}(\left|\inptext{Grid\_Yplus\_Wall}\right|)$-th
         \newline
         degree of freedom off the wall.\\
  $>0$ & This specifies the height (in \yp{}) of the first cell off the
         wall.
\end{tabular}
\end{nmldescbody}

%==========================================================================

\begin{nmldescitem}
Grid_Yplus_Of_One = 8.42e-6  ! [m]
\end{nmldescitem}
\inp*[nmlTestcaseChannel::GridYplusOfOne]{}[Grid\_Yplus\_Of\_One][]

\begin{nmldescbody}
This input specifies the dimensional length (in meters) corresponding to
\yp{}=1 to use with the internal grid generator for either of the channel flow
problems.
%
Note, for the annular-pipe channel flow case, this is applied to the radial wall
spacing.
\end{nmldescbody}

%==========================================================================

\begin{nmldescitem}
Body_Force = 0.0  ! [N]
\end{nmldescitem}
\inp*[nmlTestcaseChannel::BodyForce]{}[Body\_Force][]

\begin{nmldescbody}
Body force (in Newtons) used to drive the channel flow problem.
\end{nmldescbody}

%==========================================================================

\begin{nmldescitem}
Nondim_Output_Coord = F
\end{nmldescitem}
\inp*[nmlTestcaseChannel::NondimOutputCoord]{}[Nondim\_Output\_Coord][]

\begin{nmldescbody}
If running the channel flow problem and this input is set to true,
then all the grid coordinates written to the CGNS solution files
will be nondimensionalized by the channel half height.
\end{nmldescbody}

%==========================================================================

%\begin{nmldescitem}
%Rake_Save_Interval = 100
%\end{nmldescitem}
%\inp*[nmlTestcaseChannel::RakeSaveInterval]{}[Rake\_Save\_Interval][]
%
%\begin{nmldescbody}
%\NOTE
%(not working)
%
%This input allows the option to collect a time history of the flow
%along a rake in order to extract energy spectrums from a channel flow.
%%
%It specifies the number of time steps between saving the
%instantaneous flow profile along a rake to the time history data.
%
%\NOTE
%This can only be used with the channel flow problem using an
%internally generated grid.
%
%\vspace*{0.5em}%
%\FIXME{
%\begin{itemize}
%\item
%Where are the rake points located?
%\item
%Should this be removed for now and added again when it works?
%\end{itemize}
%}
%\end{nmldescbody}

%==========================================================================

